\section{Warp Ensemble Selection}

\subsection{Motivation}

The V3 decomposition in Section~4.4 identifies encoder Euclidean geometry as the primary bottleneck inside CEM pools. Replacing predicted terminal latents with actual terminal-image encodings improves $\Rpool$, showing that rollout error contributes, but the resulting V3 ranker remains far below the physical V1 oracle on the same candidates. Thus the pool contains rankable physical structure that Euclidean distance in the learned latent space does not expose.

We also tested thirteen inference-time interventions that targeted search or selection without pool-level learning, including learned endpoint cost heads, frozen generic encoders, hybrid oracle CEM, Subspace-CEM, diversity-regularized CEM variants, multi-view random projections, directional costs, local PCA reweighting, and earlier warp/voting variants. These interventions did not simultaneously improve pool ranking, rank-1 success, and selection regret under the locked criteria. The negative pattern motivates a calibrated selector trained directly on the object that fails: ranked candidates inside CEM-produced pools.

\subsection{Method}

Warp Ensemble Selection learns a small residual transformation of the 192-dimensional LeWM latent and then retains the Euclidean form of the terminal cost in the transformed space. Each local warp is
\begin{equation}
\mathrm{warp}(z) = z + \alpha \cdot \mathrm{MLP}_{\theta}(z),
\qquad \alpha = 0.05 .
\label{eq:warp}
\end{equation}
The MLP has one hidden layer with 32 units, ReLU activation, dropout 0.1, and a zero-initialized final layer, so training starts from the identity map. This residual parameterization keeps the calibration local rather than replacing the planner's objective with an unconstrained scalar cost.

We train ten warps using ten-fold splits over the 100 PushT CEM pools. For each fold, the warp is trained on the other 90 pools for 300 epochs using pairwise ranking supervision from simulator-scored $C_\mathrm{real\_state}$ labels. The ranking loss samples hard pairs from the top-50 candidates under $\Cmodel$, focusing training on the region where final action selection is decided. An identity regularizer with weight $3\times10^{-2}$ penalizes large deviations from the original latent geometry.

At inference time, CEM search is unchanged: proposal updates and the final 300-candidate pool are produced by the default Euclidean $\Cmodel$. Only the final rank-1 selection changes. For each candidate $i$ and warp $k$, we compute the warped Euclidean cost
\[
C_{k,i} = \lVert \mathrm{warp}_k(\hat z_{H,i}) - \mathrm{warp}_k(z_g) \rVert_2^2 .
\]
The candidates are ranked within each warp, and the selected action is the candidate with the lowest average rank across the ten warps:
\[
i^\star = \arg\min_i \frac{1}{10}\sum_{k=1}^{10} \mathrm{rank}_k(i).
\]
Thus the method changes only the final selector, not the search distribution or simulator budget.

\subsection{Results}

Table~\ref{tab:warp-selection-main} compares default CEM, Warp Ensemble Selection, and the MPPI sanity-check baseline on the same 30 PushT pairs with three seeds. The Warp Ensemble uses the default CEM pools, but selects rank-1 actions by the rank-averaged warped cost. It improves planning success by 3.3 percentage points, reduces mean selection regret by 14.6\%, and increases the pool-ranking correlation of the active selector by 0.161 (95\% paired-bootstrap CI $[0.109,0.222]$). The regret reduction is $-2.79$ in absolute cost units with 95\% CI $[-5.67,-0.60]$.

\begin{table}[t]
\centering
\small
\begin{tabular}{lccc}
\toprule
Metric & Default CEM & Warp Ensemble & MPPI \\
\midrule
Planning success (30p, 3s) & 56.7\% & \textbf{60.0\%} & 38.9\% \\
Selection regret & 19.17 & \textbf{16.37} & 31.46 \\
$\Rpool(\mathrm{selector})$ & 0.085 & \textbf{0.246} & 0.194 \\
\bottomrule
\end{tabular}
\caption{Matched PushT comparison for Warp Ensemble Selection. Default CEM and Warp Ensemble use identical CEM search; they differ only in final rank-1 selection. $\Rpool(\mathrm{selector})$ is computed from the score used by the final selector. MPPI is included as an optimizer-dependence sanity check.}
\label{tab:warp-selection-main}
\end{table}

Table~\ref{tab:warp-selection-subsets} breaks the same 30-pair evaluation down by subset. The selector does not regress on any subset. Gains are largest for V1-favorable pairs, where success rises from 26.7\% to 46.7\% and regret drops from 34.4 to 27.2. Invisible-quadrant pairs show lower regret but unchanged success, matching the earlier diagnosis that these pools often contain little or no success mass.

\begin{table}[t]
\centering
\small
\begin{tabular}{lcccc}
\toprule
Subset & Default Success & Warp Success & Default Regret & Warp Regret \\
\midrule
Invisible quadrant & 8.3\% & 8.3\% & 29.1 & \textbf{24.4} \\
Ordinary & 83.3\% & 83.3\% & 12.8 & \textbf{11.9} \\
Latent favorable & 100.0\% & 100.0\% & \textbf{3.4} & 3.5 \\
V1 favorable & 26.7\% & \textbf{46.7\%} & 34.4 & \textbf{27.2} \\
\bottomrule
\end{tabular}
\caption{Per-subset results for default CEM and Warp Ensemble Selection. Each row averages three seeds per pair before averaging over pairs in the subset.}
\label{tab:warp-selection-subsets}
\end{table}

\subsection{Discussion}

Warp Ensemble Selection validates the diagnostic interpretation of the audit. Pool-level ranking labels can calibrate the local geometry that default Euclidean costs fail to exploit, while preserving the default CEM search trajectory. The method therefore occupies a different point in the repair space from learned endpoint cost heads: it trains on pool-local rankings, keeps the cost Euclidean after a residual latent transformation, and applies only at final selection.

The method has two important limitations. First, it requires offline pools with simulator-scored $C_\mathrm{real\_state}$ labels, so it is a calibration method for audited tasks rather than a privilege-free deployment monitor. Second, it does not improve CEM search itself. If the final pool contains no successful candidate, final selection can reduce regret but cannot create success mass. This explains the invisible-quadrant pattern in Table~\ref{tab:warp-selection-subsets}. Finally, the residual structure preserves global endpoint behavior: the endpoint-ranking check changes by only $+0.020$, indicating that the gain comes from local calibration rather than a large trade-off against endpoint geometry.
